\section{MDx and SHA Functions}\label{sec:3}

The MDx functions are a form of ``Dedicated hash functions'' defined
earlier in \S\ref{sec:2}. The MD4 \cite{ref21} hash
function was introduced by Ron Rivest in 1990. While weaknesses were
discovered with MD4
\cite{ref15}\cite{ref30}\cite{ref32},
MD4's significance is that it served as the starting point for the
development of a series of similar iterated hash functions. MD4-derived
hash functions include the 256-bit extension of MD4
\cite{ref23} and the Secure Hash Algorithm (SHA) designed by
NIST/NSA, later substituted by the slightly revised SHA-1
\cite{ref10}. Additionally MD4 is optimized to be highly
efficient on 32-bit word computers and is thus able to produce 128-bit
hash results far quicker than alternative methods, e.g.~via modified
block ciphers.

To counter the weaknesses discovered with MD4 Rivest introduced MD5
\cite{ref22} in 1991 to strengthen the hash function
accordingly. MD5 differs from MD4 on the following points
\cite{ref3}.

\begin{itemize}
\tightlist
\item
  A fourth round was added.
\item
  The second round function has been changed from the majority function
  \((X\land Y)\lor(X\land Z)\lor(Y\land Z)\) to the multiplexer function
  \((X\land Z)\lor(Y\land\neg Z).\)
\item
  The order in which input words are accessed in rounds 2 and 3 was
  changed.
\item
  The shift amounts in each round have been changed. None are the same
  now.
\item
  Each step now has a unique additive constant.
\item
  Each step now adds in the result of the previous step.
\end{itemize}

Where \(\land\) denotes bit-wise AND, \(\lor\) denotes bit-wise OR,
\(\oplus\) denotes bit-wise XOR, and \(\neg Z\) denotes the bit-wise
complement of \emph{Z}.

\subsection{MD5 Detail}\label{sec:3.1}\label{md5-detail}

MD5 is perhaps the most extensively covered dedicated hash function in
the literature, likely due to its popularity. Rivest details the
algorithm and provides sample C source in \cite{ref22}.
Schneier provides a survey of several of the more frequently used hash
functions in \cite{ref7}. For our purposes a subset of that
information is sufficient.

The MD5 algorithm follows the iterative structure described in
\S\ref{sec:2.2}, where the hash is computed by repeated
application of a compression (or \emph{round}) function to successive
blocks of the message. The message to be hashed is first padded to a
multiple of 512 bits and then divided into a sequence of 512-bit message
blocks. The compression function takes two inputs, a 128-bit chaining
value and a 512-bit message block, and produces as output a new 128-bit
chaining variable, which is input to the next iteration of the
compression function \cite{ref4}. In every step one of the
chaining variables is updated. A typical step operation of the round
function is

\[
A\leftarrow B+\operatorname{ROTL}_{32}
       \bigl(A+\Phi(B,C,D)+X[j]+K,s\bigr).
\]
Each 512-bit block comprises 16 words $X[0],\ldots,X[15]$ of 32 bits;
$j$ selects a word for the current step. All additions are modulo $2^{32}$,
and $\operatorname{ROTL}_{32}(u,s)$ rotates a 32-bit word left by $s$ bits.
The Boolean function $\Phi$ and constant $K$ depend on the round and step.

The initial register values, written as ordinary hexadecimal \emph{32-bit
integers}, are \cite{ref22}
\begin{equation}\label{eq:3.1}
\begin{aligned}
A&=\mathtt{0x67452301},&B&=\mathtt{0xefcdab89},\\
C&=\mathtt{0x98badcfe},&D&=\mathtt{0x10325476}.
\end{aligned}
\end{equation}
RFC 1321 also lists their bytes in little-endian order, for example
\texttt{01 23 45 67} for $A$. The original paper printed that byte sequence
as if it were an integer; the distinction is now explicit.

There are four nonlinear functions, one used in each operation (a
different one for each round).

\begin{equation}\label{eq:3.2}
\begin{aligned}\Phi_1(X,Y,Z)&=(X\land Y)\lor(\neg X\land Z)\\\Phi_2(X,Y,Z)&=(X\land Z)\lor(Y\land\neg Z)\\\Phi_3(X,Y,Z)&=X\oplus Y\oplus Z\\\Phi_4(X,Y,Z)&=Y\oplus(X\lor\neg Z)\end{aligned}
\end{equation}

These functions are designed so that if the corresponding bits of
\emph{X, Y,} and \emph{Z} are independent and unbiased, then each bit of
the result will also be independent and unbiased. The function
\(\Phi_1\) is the bit-wise conditional: If \emph{X} then \emph{Y} else
\emph{Z}. The function \(\Phi_2\) is the bit-wise conditional: If
\emph{Z} then \emph{X} else \emph{Y}. The function \(\Phi_3\) is the
bit-wise parity operator \cite{ref7}. After the last message
block has been processed, the final chaining value is output as the hash
of the message.

\subsection{SHA Functions}\label{sec:3.2}\label{sha-functions}

The Secure Hash Algorithm (SHA) \cite{ref10} is a digest
algorithm proposed by the US NIST (National Institute of Standards and
Technology) agency as a standard digest algorithm in August 2002. NIST
released a first version of this algorithm --- referred to as "SHA" back
in 1992. It subsequently discovered a weakness with SHA and released a
second version which it referred to as "SHA-1''.

Where SHA-1 has a 160 bit hash value, NIST also defined hash functions
with larger hash output such as SHA-224 (224 bits), SHA-256 (256 bits),
SHA-384 (384 bits) and SHA-512 (512 bits). The cryptographic community
generally considers SHA-256, SHA-384 and SHA-512 to be stronger than
MD5.

Detail on the internals of the SHA family can be found in
\cite{ref10}.
